% \onecolumn

\chapter{First Appendix}
\label{sec:apx:first_appendix}

This appendix collects the reference material the main text points to: exact
configurations, tool versions and results too detailed for the body chapters. Sections are
ordered by how likely each is to hold real content before submission.

\section{Synthesis Algorithm Configurations}
\label{sec:apx:algorithms}
This section records the concrete invocations
behind~\ref{sec:method:data:generation} and~\ref{sec:prelim:algorithms}. Every entry below is
transcribed from the generation pipeline itself rather than from its documentation; where
the two disagree, the pipeline is authoritative and the disagreement is noted.

\subsection*{Transformation recipes}
Each recipe is $20$ commands drawn independently and uniformly from
\texttt{balance}, \texttt{rewrite}, \texttt{rewrite -z}, \texttt{refactor},
\texttt{refactor -z}, \texttt{resub} and \texttt{resub -z}, where \texttt{-z} enables
zero-cost replacements. The $200$ recipes used are the first $200$ of the $1{,}500$
released with OpenABC-D \citep{chowdhury2021openabc} and are applied unmodified. The
generator strips the input, output and hashing commands from each released recipe and
appends a \texttt{write\_aiger} after every surviving command, which is what makes each
recipe prefix a stored graph.

Each released recipe continues past the $20$ transformations into \texttt{dch} and then a
technology-mapping tail, \texttt{map -B 0.9}, \texttt{topo}, \texttt{stime -c},
\texttt{buffer -c}, \texttt{upsize -c} and \texttt{dnsize -c}. The stripping rule does not
remove that tail, so $27$ write commands are emitted per recipe, but no standard-cell
library is loaded and the mapped network cannot be written as an AIG, so only the $21$
graphs up to and including \texttt{dch} are produced. The count assertion in the generation
job, exactly $4{,}201$ AIGs per design, is what establishes that the remaining six writes
produced nothing.

\subsection*{Optimization invocations}
Each algorithm is applied through the template in Table~\ref{tab:apx:algorithms}, with the
tool configuration file sourced first in every case.

\begin{table}[htbp]
\centering
\caption{Optimization invocations, one per algorithm, as issued by the generation pipeline.}
\label{tab:apx:algorithms}
\footnotesize
\begin{tabular}{ll}
\toprule
\textbf{Algorithm} & \textbf{Command} \\
\midrule
\texttt{Orchestrate} & \texttt{strash; orchestrate} \\
\texttt{Deepsyn}     & \texttt{strash; \&get; \&deepsyn -S <seed> -T 20; \&put} \\
\texttt{Syn4}        & \texttt{strash; \&get; \&syn4; \&put} \\
\texttt{C2RS}        & \texttt{strash; c2rs} \\
\bottomrule
\end{tabular}
\end{table}

Three of the four are invoked with no parameters at all, so their behaviour is the tool's
own default in each case. The pass budgets and scoring metrics named in the dataset
documentation are not passed on the command line and were not set by this work.
\texttt{C2RS} is an alias expanding to
\texttt{b -l; rs -K 6 -l; rw -l; rs -K 6 -N 2 -l; rf -l; rs -K 8 -l; b -l;
rs -K 8 -N 2 -l; rw -l; rs -K 10 -l; rwz -l; rs -K 10 -N 2 -l; b -l; rs -K 12 -l;
rfz -l; rs -K 12 -N 2 -l; rwz -l; b -l}, which is the source of the widening cut size
reported in~\ref{sec:prelim:algorithms:others}.

\texttt{Deepsyn} is the one stochastic algorithm. Its seed is drawn from the shell random
number generator once per input graph, is interpolated into the command, and is written
neither to the run log nor to the metadata tables, so those graphs cannot be regenerated
(\ref{sec:discussion:limitations:reproducibility}).

The \texttt{abc} binary is built from source on the cluster, not installed through a
versioned module. The caller supplies its path (\texttt{\$HOME/abc/abc} by default), and no
submodule pin, lockfile entry or recorded commit hash fixes which revision that is. No
invocation in the generation pipeline captures a version string, so the exact tool behind
every command in Table~\ref{tab:apx:algorithms} is unrecorded and unrecoverable, a standing
bound on reproducibility (\ref{sec:discussion:limitations:reproducibility}): the commands are
exact, but the version that ran them is not.
% TODO: pin the abc commit or version if it is ever recovered from the build host.

\section{Supplementary Figures}
\label{sec:apx:figures}
% TODO: figures and tables moved from the results chapter as diagnostic rather than load-bearing, starting with the hyperparameter sweep and sensitivity figures.

\section{Experimental Configuration}
\label{sec:apx:config}
% TODO: full run configuration table: seeds, split settings, batch budget, optimizer settings and precision.

\section{Hyperparameter Search Space}
\label{sec:apx:hpspace}
% TODO: full table of the swept hyperparameter dimensions and ranges.

\section{Library and Tool Versions}
\label{sec:apx:versions}
% TODO: table of every library and tool version behind a reported result: Python toolchain, GNN framework, aigverse, mockturtle, Optuna.

\section{Reduction Method Parameters}
\label{sec:apx:params}
% TODO: table of every reduction method's parameters and whether each was calibrated against measured retention or left at a default.

\section{Baseline Hyperparameters}
\label{sec:apx:baselines}
% TODO: SynthNet, HOGA and DeepGate4 hyperparameters as used versus each paper's published values, with the reason for every deviation (src/train_baseline.py, src/baselines/*/regressor.py).

\section{Extended Results}
\label{sec:apx:results}
% TODO: per-configuration tables too detailed for the results chapter: full metric sets per method, per-tier and per-design error breakdowns, and offline retention per method and shard.
